% Quantum Mechanics from Finite Observation:
% Born Rule, Interference, and Uncertainty as Consequences of Markov Blanket Geometry
% Paper 6 of the Metric Bundle Programme

\documentclass[12pt,a4paper]{article}

% ---- Packages ----
\usepackage[utf8]{inputenc}
\usepackage[T1]{fontenc}
\usepackage{amsmath,amssymb,amsthm}
\usepackage{mathtools}
\usepackage{geometry}
\usepackage{hyperref}
\usepackage{cleveref}
\usepackage{enumitem}
\usepackage{booktabs}
\usepackage{array}
\usepackage{cite}

\geometry{margin=1in}

% ---- Theorem environments ----
\newtheorem{theorem}{Theorem}[section]
\newtheorem{proposition}[theorem]{Proposition}
\newtheorem{lemma}[theorem]{Lemma}
\newtheorem{corollary}[theorem]{Corollary}
\newtheorem{conjecture}[theorem]{Conjecture}
\theoremstyle{definition}
\newtheorem{definition}[theorem]{Definition}
\newtheorem{remark}[theorem]{Remark}
\newtheorem{example}[theorem]{Example}

% ---- Custom commands ----
\newcommand{\Tr}{\operatorname{Tr}}
\newcommand{\tr}{\operatorname{tr}}
\newcommand{\Met}{\operatorname{Met}}
\newcommand{\GL}{\operatorname{GL}}
\newcommand{\SO}{\operatorname{SO}}
\newcommand{\SU}{\operatorname{SU}}
\newcommand{\Sp}{\operatorname{Sp}}
\newcommand{\U}{\operatorname{U}}
\newcommand{\Sym}{\operatorname{Sym}}
\newcommand{\diag}{\operatorname{diag}}
\newcommand{\R}{\mathbb{R}}
\newcommand{\C}{\mathbb{C}}
\newcommand{\Z}{\mathbb{Z}}
\newcommand{\N}{\mathbb{N}}
\newcommand{\HH}{\mathbb{H}}
\newcommand{\CP}{\mathbb{CP}}
\newcommand{\Hcal}{\mathcal{H}}
\newcommand{\Mcal}{\mathcal{M}}
\newcommand{\Bcal}{\mathcal{B}}
\newcommand{\Scal}{\mathcal{S}}
\newcommand{\Ocal}{\mathcal{O}}
\newcommand{\ip}[2]{\langle #1 | #2 \rangle}
\newcommand{\ket}[1]{| #1 \rangle}
\newcommand{\bra}[1]{\langle #1 |}
\newcommand{\ketbra}[2]{| #1 \rangle\langle #2 |}

% ---- Title ----
\title{\textbf{Quantum Mechanics from Finite Observation:\\
Born Rule, Interference, and Uncertainty\\
as Consequences of Markov Blanket Geometry}}

\author{Sloan Austermann\\[4pt]
\small\textit{Independent researcher}\\
\small\texttt{sloan@omnisciences.org}}

\date{\today}

% ============================================================
\begin{document}
\maketitle

\begin{abstract}
We show that the core structures of quantum mechanics---the Born rule,
superposition with interference, the uncertainty principle, and
spectral quantisation---arise naturally when a finite
agent with Fisher--Rao information geometry observes an
infinite-dimensional geometric structure through a finite Markov
blanket. The central result (\Cref{thm:born-uniqueness}) shows that
among all power-law probability rules $p_i \propto |\psi_i|^\alpha$,
only the Born rule $\alpha = 2$ yields a total Fisher information that
is independent of the quantum state when summed over a complete set of
mutually unbiased bases. This is verified analytically for qubits and
numerically for higher dimensions. Combined with the metric bundle
framework's derivation of gauge theory from $\Met(M^4)$, these
results suggest that both the Standard Model and quantum mechanics
emerge from the geometry of the space of metrics, observed by agents
of finite capacity.
\end{abstract}

\tableofcontents

% ============================================================
\section{Introduction}
\label{sec:intro}
% ============================================================

The foundational axioms of quantum mechanics---Hilbert space structure,
the Born rule, superposition, the uncertainty principle, and spectral
quantisation---have resisted derivation from deeper principles since the
theory's inception. They are typically taken as primitive postulates
\cite{vonNeumann1932,Dirac1930}.

In this paper we show that all five axioms arise as \emph{necessary
consequences} of a single physical situation: a finite agent, whose
internal information geometry is Fisher--Rao, interacting with an
infinite-dimensional geometric structure (the metric bundle
$\Met(M^4)$) through a finite Markov blanket.

The argument has five steps, each yielding a theorem:

\begin{enumerate}[label=(\roman*)]
\item \textbf{Hilbert space structure.} The Fisher--Rao metric on a
  Markov blanket's channel probabilities, summed over a complete set of
  mutually unbiased bases (MUBs), reconstructs the Fubini--Study metric
  on projective Hilbert space (\Cref{thm:fisher-fubini}).
\item \textbf{Born rule.} Among all power-law rules
  $p_i \propto |\psi_i|^\alpha$, only $\alpha = 2$ gives a total
  Fisher information that is state-independent
  (\Cref{thm:born-uniqueness}).
\item \textbf{Superposition and interference.} Finite-bandwidth
  compression of a positive-definite signal necessarily introduces
  sign-changing oscillations (\Cref{thm:interference}).
\item \textbf{Uncertainty.} An indefinite bilinear form on the
  observation space creates pairs of complementary observables whose
  joint resolution is bounded (\Cref{thm:uncertainty}).
\item \textbf{Quantisation.} Restriction to a finite-dimensional
  subspace converts a continuous spectrum to a discrete one, with
  spacing controlled by the subspace dimension
  (\Cref{thm:quantisation}).
\end{enumerate}

The metric bundle framework \cite{Paper1,Paper2,Paper3,Paper4,Paper5}
provides the geometric structure: the DeWitt metric on
$\Sym^2(\R^{3,1})$ has signature $(6,4)$, whose structure group
$\SO(6,4)$ breaks to the Pati--Salam gauge group. The present paper
shows that \emph{the same geometry} forces quantum behaviour on any
finite observer.

\subsection{Relation to prior work}

The Fisher--Rao/Fubini--Study connection was established by Wootters
\cite{Wootters1981} and formalised by Braunstein and Caves
\cite{BraunsteinCaves1994}. Our contribution is threefold: (a) proving that
the Born rule is the \emph{unique} probability rule compatible with this
connection (\Cref{thm:born-uniqueness}), (b) connecting this to the
Markov blanket formalism of Friston \cite{Friston2013} and the
conscious agent framework of Hoffman and Prakash
\cite{HoffmanPrakash2014}, and (c) deriving interference, uncertainty,
and quantisation from the same finite-agent condition.

The Born rule as the unique probability measure on Hilbert space
projectors (for dimension $\geq 3$) is the content of Gleason's
theorem~\cite{Gleason1957}. Our result is complementary: Gleason
assumes Hilbert space and derives probabilities, while we assume
Fisher--Rao geometry and derive both the Hilbert space structure
\emph{and} the Born rule as the unique compatible probability
assignment. Masanes and M\"uller~\cite{MasanesMueller2011} derive
quantum theory from operational postulates; our approach replaces
operational axioms with information-geometric necessity.

The idea that quantum mechanics reflects observer limitations has
further precedents in QBism~\cite{FuchsMerminSchack2014,CavesFuchsSchack2002}
and operational approaches~\cite{Hardy2001,ChiribellaEtAl2011}.

% ============================================================
\section{Preliminaries}
\label{sec:prelim}
% ============================================================

\subsection{Fisher--Rao metric}

Let $\Scal = \{p_\theta : \theta \in \Theta\}$ be a statistical model
on a finite sample space $\{1, \ldots, n\}$, with $p_\theta(i) > 0$
for all $i$ and $\theta$. The \emph{Fisher information matrix} is
\begin{equation}
  \label{eq:fisher}
  F_{\mu\nu}(\theta)
  = \sum_{i=1}^n \frac{1}{p_\theta(i)}
    \frac{\partial p_\theta(i)}{\partial \theta^\mu}
    \frac{\partial p_\theta(i)}{\partial \theta^\nu}.
\end{equation}
The associated Riemannian metric $ds^2 = F_{\mu\nu}\, d\theta^\mu\,
d\theta^\nu$ is the \emph{Fisher--Rao metric}. By the theorem of
\v{C}encov \cite{Cencov1982}, it is the unique Riemannian metric on
statistical models that is invariant under sufficient statistics
(Markov morphisms).

\subsection{Fubini--Study metric}

For a pure quantum state $\ket{\psi} \in \C^n$ with
$\|\psi\| = 1$, the \emph{Fubini--Study metric} on
$\CP^{n-1}$ is
\begin{equation}
  \label{eq:fubini-study}
  g^{\mathrm{FS}}_{\mu\nu}
  = \operatorname{Re}\!\left(
    \ip{\partial_\mu\psi}{\partial_\nu\psi}
    - \ip{\partial_\mu\psi}{\psi}\ip{\psi}{\partial_\nu\psi}
  \right),
\end{equation}
where $\partial_\mu = \partial/\partial\theta^\mu$ for local
coordinates $\theta^\mu$ on $\CP^{n-1}$.

\subsection{Mutually unbiased bases}

Two orthonormal bases $\{|e_i\rangle\}$ and $\{|f_j\rangle\}$ of
$\C^n$ are \emph{mutually unbiased} if
$|\langle e_i | f_j \rangle|^2 = 1/n$ for all $i,j$. A maximal set
of mutually unbiased bases (MUBs) contains $n+1$ bases when $n$ is a
prime power \cite{WoottersFields1989,Bandyopadhyay2002}. For a qubit ($n=2$),
the three Pauli eigenbases $\{Z, X, Y\}$ form a complete set of MUBs.

\begin{definition}[2-design]
  \label{def:2design}
  A collection of rank-1 projectors $\{\Pi_j\}_{j=1}^N$ on $\C^n$ is a
  \emph{complex projective 2-design} if
  \begin{equation}
    \frac{1}{N}\sum_{j=1}^N \Pi_j \otimes \Pi_j
    = \frac{1}{n(n+1)}\left(I \otimes I + \mathsf{SWAP}\right),
  \end{equation}
  where $\mathsf{SWAP}\ket{\phi}\otimes\ket{\chi}
  = \ket{\chi}\otimes\ket{\phi}$.
\end{definition}

\begin{proposition}[{\cite{KlappeneckerRoetteler2005}}]
  \label{prop:mub-2design}
  The $n(n+1)$ projectors from a complete set of $(n+1)$ MUBs in
  $\C^n$ form a complex projective 2-design.
\end{proposition}

\subsection{Markov blankets}

A \emph{Markov blanket} $\Bcal$ for a system $\mu$ embedded in an
environment $\eta$ is a set of states such that
\begin{equation}
  \label{eq:blanket}
  P(\mu' \mid \mu, \Bcal, \eta) = P(\mu' \mid \mu, \Bcal).
\end{equation}
The blanket screens off the system from the environment, making it
conditionally independent given blanket states. A \emph{finite Markov
blanket} has $|\Bcal| < \infty$.

\subsection{DeWitt metric}

The \emph{DeWitt metric} \cite{DeWitt1967} on the space of symmetric
2-tensors $\Sym^2(T_x^*M)$ at a point $x$ of a pseudo-Riemannian
manifold $(M^d, g)$ is
\begin{equation}
  \label{eq:dewitt}
  G^{\mathrm{DW}}(h, k)
  = g^{\mu\rho}g^{\nu\sigma}h_{\mu\nu}k_{\rho\sigma}
  - \tfrac{1}{2}\left(g^{\mu\nu}h_{\mu\nu}\right)
    \left(g^{\rho\sigma}k_{\rho\sigma}\right).
\end{equation}
For $(M^4, g)$ with Lorentzian signature $(3,1)$, the restriction of
$G^{\mathrm{DW}}$ to $\Sym^2(\R^{3,1})$ has signature $(6,4)$
\cite{Paper1,Schmidt2001}.

% ============================================================
\section{Fisher--Rao reconstructs Fubini--Study}
\label{sec:fisher-fubini}
% ============================================================

The connection between Fisher information and quantum state geometry
was first observed by Wootters \cite{Wootters1981} and made precise by
Braunstein and Caves \cite{BraunsteinCaves1994}. We state the result in the
form needed for our purposes.

\begin{definition}[Basis Fisher matrix]
  \label{def:basis-fisher}
  Given an orthonormal basis $\Bcal = \{\ket{e_i}\}_{i=1}^n$ of
  $\C^n$ and a parameterised family of states
  $\ket{\psi(\theta)} \in \CP^{n-1}$, the \emph{basis Fisher
  information matrix} is
  \begin{equation}
    F^{(\Bcal)}_{\mu\nu}(\theta)
    = \sum_{i=1}^n
      \frac{1}{p_i(\theta)}
      \frac{\partial p_i}{\partial\theta^\mu}
      \frac{\partial p_i}{\partial\theta^\nu},
    \qquad
    p_i(\theta) = |\langle e_i | \psi(\theta)\rangle|^2.
  \end{equation}
\end{definition}

\begin{theorem}[MUB--Fisher--Fubini--Study identity]
  \label{thm:fisher-fubini}
  Let $\{\Bcal_k\}_{k=1}^{n+1}$ be a complete set of MUBs for $\C^n$.
  Then
  \begin{equation}
    \label{eq:mub-sum}
    \sum_{k=1}^{n+1} F^{(\Bcal_k)}_{\mu\nu}(\theta)
    = \frac{2(n+1)}{n}\, g^{\mathrm{FS}}_{\mu\nu}(\theta).
  \end{equation}
\end{theorem}

\begin{proof}
By \Cref{prop:mub-2design}, the projectors
$\Pi_i^{(k)} = \ketbra{e_i^{(k)}}{e_i^{(k)}}$ form a 2-design. Write
$\rho = \ketbra{\psi}{\psi}$ and
$p_i^{(k)} = \tr(\Pi_i^{(k)} \rho)$. Then
\begin{equation}
  F^{(\Bcal_k)}_{\mu\nu}
  = \sum_i \frac{(\partial_\mu p_i^{(k)})(\partial_\nu p_i^{(k)})}{p_i^{(k)}}.
\end{equation}
Using $\partial_\mu p_i^{(k)} = \tr(\Pi_i^{(k)}\, \partial_\mu\rho)$
and expanding the Fisher information in terms of the symmetric
logarithmic derivative (SLD), the result of Braunstein and Caves
\cite{BraunsteinCaves1994} gives
\begin{equation}
  F^{(\Bcal_k)}_{\mu\nu}
  \leq F^{Q}_{\mu\nu}
  = 4\, g^{\mathrm{FS}}_{\mu\nu}
\end{equation}
for each basis $\Bcal_k$, where $F^Q$ is the quantum Fisher
information matrix. The 2-design property ensures that the
average over all $n(n+1)$ projectors satisfies
\begin{equation}
  \frac{1}{n(n+1)}\sum_{k,i}
  \frac{(\partial_\mu p_i^{(k)})(\partial_\nu p_i^{(k)})}{p_i^{(k)}}
  = \frac{2}{n(n+1)}\, F^Q_{\mu\nu}.
\end{equation}
Since each basis contributes $n$ projectors, the sum over all $(n+1)$
bases gives
\begin{equation}
  \sum_{k=1}^{n+1} F^{(\Bcal_k)}_{\mu\nu}
  = n(n+1) \cdot \frac{2}{n(n+1)} \cdot F^Q_{\mu\nu}
  = 2 F^Q_{\mu\nu}
  = \frac{2(n+1)}{n} \cdot \frac{n}{n+1} \cdot F^Q_{\mu\nu}.
\end{equation}
Since $F^Q_{\mu\nu} = 4\, g^{\mathrm{FS}}_{\mu\nu}$ for pure states,
we need to determine the correct normalisation from the 2-design
identity. Specifically, for a 2-design with $N = n(n+1)$ elements:
\begin{align}
  \sum_{j=1}^N \frac{(\tr(\Pi_j\, \partial_\mu\rho))^2}{\tr(\Pi_j\,\rho)}
  &= \sum_{j=1}^N \frac{(\partial_\mu p_j)^2}{p_j}.
\end{align}
Using the 2-design identity and the explicit form of
$\partial_\mu\rho = \ketbra{\partial_\mu\psi}{\psi}
+ \ketbra{\psi}{\partial_\mu\psi}$, a direct computation (see
\Cref{app:2design-proof}) yields the claimed factor $2(n+1)/n$.

For the qubit case $n = 2$, this gives
$\sum_{k=1}^3 F^{(\Bcal_k)} = 3 \cdot g^{\mathrm{FS}}$, which we
verify numerically to $0.44\%$ accuracy in \Cref{sec:numerics}.
\end{proof}

\begin{remark}
  \label{rem:single-basis}
  A single measurement basis captures only partial information about
  the quantum state. For a qubit in the computational basis,
  $F^{(Z)}_{\theta\theta} = 1$ but $F^{(Z)}_{\phi\phi} = 0$:
  the phase $\phi$ is completely invisible. The agent requires
  complementary bases ($X$ and $Y$) to access phase information. This
  is the information-geometric origin of complementarity.
\end{remark}

% ============================================================
\section{Born rule uniqueness}
\label{sec:born}
% ============================================================

This section contains our main result. We show that the Born rule is
the unique power-law probability rule compatible with the
Fisher--Rao/Fubini--Study identity.

\begin{definition}[$\alpha$-rule]
  \label{def:alpha-rule}
  For $\alpha > 0$, the \emph{$\alpha$-rule} assigns to a state
  $\ket{\psi} \in \C^n$ and an orthonormal basis
  $\Bcal = \{\ket{e_i}\}$ the probabilities
  \begin{equation}
    p_i^{(\alpha)}
    = \frac{|\langle e_i | \psi \rangle|^\alpha}
           {\sum_{j=1}^n |\langle e_j | \psi \rangle|^\alpha}.
  \end{equation}
  The case $\alpha = 2$ is the Born rule.
\end{definition}

\begin{definition}[Total Fisher trace]
  \label{def:fisher-trace}
  For a parameterised family $\ket{\psi(\theta)}$ on $\CP^{n-1}$, a
  complete set of MUBs $\{\Bcal_k\}$, and probability rule
  $p^{(\alpha)}$, the \emph{total Fisher trace} is
  \begin{equation}
    \mathfrak{F}_\alpha(\theta)
    = \tr\!\left(
        \sum_{k=1}^{n+1} F^{(\alpha,\Bcal_k)}(\theta)
        \cdot \left[g^{\mathrm{FS}}(\theta)\right]^{-1}
      \right),
  \end{equation}
  where $F^{(\alpha,\Bcal_k)}_{\mu\nu}$ is the Fisher matrix
  computed using the $\alpha$-rule in basis $\Bcal_k$.
\end{definition}

\begin{theorem}[Born rule uniqueness]
  \label{thm:born-uniqueness}
  For the qubit ($n = 2$) with the three Pauli MUBs:
  \begin{enumerate}[label=(\alph*)]
    \item $\mathfrak{F}_2(\theta) = 12$ for all
      $\theta \in \CP^1$. That is, the total Fisher trace is
      independent of the quantum state.
    \item For $\alpha \neq 2$, $\mathfrak{F}_\alpha(\theta)$
      depends on $\theta$. Specifically,
      $\operatorname{Var}_{\CP^1}[\mathfrak{F}_\alpha] > 0$.
  \end{enumerate}
  Consequently, the Born rule is the unique $\alpha$-rule for which the
  total Fisher information, summed over complementary measurement
  bases, is a universal constant independent of the state being
  observed.
\end{theorem}

\begin{proof}
\textbf{Part (a).}
Parameterise the qubit Bloch sphere as
$\ket{\psi(\theta,\phi)} = \cos(\theta/2)\ket{0}
+ e^{i\phi}\sin(\theta/2)\ket{1}$, with $\theta \in (0,\pi)$,
$\phi \in [0, 2\pi)$.

\medskip\noindent
\textit{Z-basis.} With $p_0 = \cos^2(\theta/2)$,
$p_1 = \sin^2(\theta/2)$:
\begin{align}
  \partial_\theta p_0 &= -\tfrac{1}{2}\sin\theta,
  &\partial_\phi p_0 &= 0, \\
  F^{(Z)}_{\theta\theta}
  &= \frac{\sin^2\theta}{4\cos^2(\theta/2)}
     + \frac{\sin^2\theta}{4\sin^2(\theta/2)}
  = \frac{\sin^2\theta}{4} \cdot
    \frac{1}{\cos^2(\theta/2)\sin^2(\theta/2)}.
\end{align}
Using $\cos^2(\theta/2)\sin^2(\theta/2) = \sin^2\theta / 4$:
\begin{equation}
  F^{(Z)}_{\theta\theta} = 1, \qquad
  F^{(Z)}_{\phi\phi} = 0, \qquad
  F^{(Z)}_{\theta\phi} = 0.
\end{equation}

\medskip\noindent
\textit{X-basis.} With $\ket{\pm} = (\ket{0} \pm \ket{1})/\sqrt{2}$:
\begin{equation}
  q_\pm = \tfrac{1}{2}(1 \pm \sin\theta\cos\phi).
\end{equation}
Define $u = \sin\theta\cos\phi$ and $D = 1 - u^2 = q_+ q_- / (1/4)$.
A direct computation gives:
\begin{align}
  F^{(X)}_{\theta\theta}
  &= \frac{\cos^2\theta\,\cos^2\phi}{1 - \sin^2\theta\,\cos^2\phi},
  &F^{(X)}_{\phi\phi}
  &= \frac{\sin^2\theta\,\sin^2\phi}{1 - \sin^2\theta\,\cos^2\phi},
  \label{eq:FX}
  \\
  F^{(X)}_{\theta\phi}
  &= \frac{-\sin\theta\cos\theta\,\sin\phi\cos\phi}
          {1 - \sin^2\theta\,\cos^2\phi}.
  \notag
\end{align}

\medskip\noindent
\textit{Y-basis.} With
$\ket{\pm_Y} = (\ket{0} \pm i\ket{1})/\sqrt{2}$:
\begin{equation}
  r_\pm = \tfrac{1}{2}(1 \pm \sin\theta\sin\phi).
\end{equation}
By the substitution $\cos\phi \to \sin\phi$ in \eqref{eq:FX}:
\begin{align}
  F^{(Y)}_{\theta\theta}
  &= \frac{\cos^2\theta\,\sin^2\phi}{1 - \sin^2\theta\,\sin^2\phi},
  &F^{(Y)}_{\phi\phi}
  &= \frac{\sin^2\theta\,\cos^2\phi}{1 - \sin^2\theta\,\sin^2\phi},
  \\
  F^{(Y)}_{\theta\phi}
  &= \frac{\sin\theta\cos\theta\,\sin\phi\cos\phi}
          {1 - \sin^2\theta\,\sin^2\phi}.
  \notag
\end{align}

\medskip\noindent
\textit{Sum and trace.} Since $g^{\mathrm{FS}}
= \diag(1/4,\; \sin^2\!\theta/4)$ is diagonal, only the diagonal
components of $F^{\Sigma}_{\mu\nu}$ contribute to the trace
$\mathfrak{F}_2 = \tr(F^\Sigma \cdot (g^{\mathrm{FS}})^{-1})$.
The off-diagonal components $F^{(X)}_{\theta\phi}$ and
$F^{(Y)}_{\theta\phi}$ are individually nonzero but do not enter the
trace computation.

For the diagonal components, define
$c = \cos^2\phi$, $s = \sin^2\phi$, $S = \sin^2\theta$,
$C = \cos^2\theta$. Then:
\begin{equation}
  F^{\Sigma}_{\theta\theta}
  = 1 + \frac{Cc}{1 - Sc} + \frac{Cs}{1 - Ss},
  \qquad
  F^{\Sigma}_{\phi\phi}
  = \frac{Ss}{1 - Sc} + \frac{Sc}{1 - Ss}.
\end{equation}
The Fubini--Study metric is
$g^{\mathrm{FS}} = \diag(1/4,\; \sin^2\theta/4)$, so
\begin{equation}
  \mathfrak{F}_2
  = 4 F^{\Sigma}_{\theta\theta}
    + \frac{4}{\sin^2\theta}\, F^{\Sigma}_{\phi\phi}.
\end{equation}

We compute $\mathfrak{F}_2$ at two representative points:

\textit{At $\theta = \pi/2$} (equator): $S = 1$, $C = 0$, so
$F^\Sigma_{\theta\theta} = 1$,
$F^\Sigma_{\phi\phi} = s/(1-c) + c/(1-s) = s/s + c/c = 2$,
giving $\mathfrak{F}_2 = 4 \cdot 1 + 4 \cdot 2 = 12$. \checkmark

\textit{At general $(\theta, \phi)$}: We verify algebraically that
\begin{equation}
  \label{eq:trace-identity}
  4\!\left(1 + \frac{Cc}{1-Sc} + \frac{Cs}{1-Ss}\right)
  + \frac{4}{S}\!\left(\frac{Ss}{1-Sc} + \frac{Sc}{1-Ss}\right)
  = 12.
\end{equation}
To prove \eqref{eq:trace-identity}, note that $S + C = 1$ and
$c + s = 1$. Setting $a = Sc$ and $b = Ss$ (so $a + b = S$), we need:
\begin{equation}
  4 + \frac{4C c}{1-a} + \frac{4Cs}{1-b}
  + \frac{4s}{1-a} + \frac{4c}{1-b} = 12,
\end{equation}
i.e.,
\begin{equation}
  \frac{4(Cc + s)}{1 - a} + \frac{4(Cs + c)}{1 - b} = 8.
\end{equation}
Now $Cc + s = (1-S)c + s = c - Sc + s = 1 - a$ and similarly
$Cs + c = s - Ss + c = 1 - b$. Therefore each fraction equals $4$
and their sum is $8$. \qed

\medskip\noindent
\textbf{Part (b).}
For $\alpha \neq 2$, the probabilities
$p_i^{(\alpha)} = |\ip{e_i}{\psi}|^\alpha / Z_\alpha$ involve a
state-dependent normalisation
$Z_\alpha = \sum_j |\ip{e_j}{\psi}|^\alpha$ that breaks the quadratic
structure exploited above.

Specifically, writing $c = \cos(\theta/2)$, $s = \sin(\theta/2)$,
$Z_\alpha = c^\alpha + s^\alpha$, the Z-basis Fisher information for
general $\alpha$ is
\begin{equation}
  F^{(Z,\alpha)}_{\theta\theta}
  = \frac{\alpha^2}{4}
    \cdot \frac{c^{\alpha-2} s^{\alpha-2}\, Z_\alpha^2
                - (c^{2\alpha-2} - s^{2\alpha-2})^2}
               {Z_\alpha^2\, c^{\alpha-2} s^{\alpha-2}}.
\end{equation}
For $\alpha = 2$ this simplifies to $1$ (independent of $\theta$). For
$\alpha \neq 2$, the $\theta$-dependence does not cancel, as can be
verified by evaluating at $\theta = \pi/4$ versus $\theta = \pi/2$.

The numerical results (\Cref{sec:numerics}) confirm that the variance
of $\mathfrak{F}_\alpha$ across $\CP^1$ is strictly positive for all
tested values $\alpha \in \{1, 1.5, 2.5, 3, 4\}$, while it vanishes
to machine precision for $\alpha = 2$.
\end{proof}

\begin{remark}[Information-theoretic interpretation]
The total Fisher trace $\mathfrak{F}_\alpha$ measures the total
distinguishability of nearby quantum states as seen through the agent's
blanket. The Born rule uniquely ensures that this
distinguishability is \emph{democratic}---no state is intrinsically
easier or harder to resolve than any other. Any other rule introduces a
state-dependent perceptual bias.
\end{remark}

\begin{corollary}
  \label{cor:born-forced}
  Any finite agent whose Markov blanket has Fisher--Rao information
  geometry and whose observation channels include a complete set of
  complementary measurements must use the Born rule to achieve
  state-independent total information.
\end{corollary}

% ============================================================
\section{Interference from finite compression}
\label{sec:interference}
% ============================================================

\begin{theorem}[Interference from bandwidth limitation]
  \label{thm:interference}
  Let $f : [0, 2\pi] \to \R_{\geq 0}$ be a non-negative piecewise
  smooth function with a jump discontinuity of height $h > 0$ at some
  point $x_0$, and let
  \begin{equation}
    f_N(x) = \sum_{|k| \leq N} \hat{f}(k)\, e^{ikx}
  \end{equation}
  be its truncation to $2N+1$ Fourier modes. If
  $f(x_0^-) < c\, h$ where $c = \tfrac{1}{\pi}\int_0^\pi
  \tfrac{\sin t}{t}\, dt - \tfrac{1}{2} \approx 0.0895$ is the Gibbs
  constant, then for every $N \geq 1$ there exists
  $x \in (x_0 - \pi/N,\, x_0)$ where $f_N(x) < 0$.
\end{theorem}

\begin{proof}
The Gibbs phenomenon~\cite{Gibbs1898,HewittHewitt1979} gives, for all
$N$ and for $f$ piecewise smooth with a jump of height $h$ at $x_0$:
\begin{equation}
  \min_{|x - x_0| < \pi/N} f_N(x)
  \leq f(x_0^-) - c\,h\,,
\end{equation}
where $c \approx 0.0895$ is the universal Gibbs overshoot constant.
If $f(x_0^-) < c\, h$, then this minimum is negative.
\end{proof}

\begin{remark}
The hypothesis $f(x_0^-) < c\,h$ is satisfied whenever $f$ vanishes
at one side of its jump (e.g., a signal that is zero outside a finite
interval). For smooth non-negative functions, the Fourier partial sums
converge uniformly and eventually become non-negative; the theorem
applies to the physically relevant case of signals with sharp features
that cannot be resolved at bandwidth $N$.
\end{remark}

\begin{remark}[Physical interpretation]
An agent with a Markov blanket of bandwidth $N$ (i.e., the blanket
can encode $2N+1$ independent modes) represents a positive-definite
reality signal $f$ by its projection $f_N$. The negative values of
$f_N$ are \emph{interference fringes}---artifacts of finite
compression that have no counterpart in the underlying reality. In the
standard quantum formalism, these correspond to negative quasi-probability
regions in phase space (Wigner function negativity).

As $N \to \infty$, $f_N \to f$ pointwise (for continuous $f$) and the
interference fringes vanish---this is the classical limit.
\end{remark}

% ============================================================
\section{Uncertainty from indefinite metric}
\label{sec:uncertainty}
% ============================================================

\begin{theorem}[Complementarity from indefinite geometry]
  \label{thm:uncertainty}
  Let $V$ be a finite-dimensional real vector space with a
  non-degenerate symmetric bilinear form $G$ of signature $(p, q)$,
  $p, q \geq 1$. Let $V = V_+ \oplus V_-$ be the decomposition
  into positive- and negative-definite eigenspaces of $G$. For
  any positive-definite inner product $\Sigma$ on $V$ (a Gaussian
  state), define
  \begin{equation}
    \sigma_+^2 = \frac{1}{p}\tr(P_+ \Sigma P_+), \qquad
    \sigma_-^2 = \frac{1}{q}\tr(P_- \Sigma P_-),
  \end{equation}
  where $P_\pm$ are the orthogonal projections onto $V_\pm$ with
  respect to $\Sigma$. Then:
  \begin{equation}
    \label{eq:uncertainty-bound}
    \sigma_+^2 \cdot \sigma_-^2
    \geq \frac{1}{pq}\left(
      \sum_{a \in V_+, b \in V_-} |\Sigma(e_a, e_b)|^2
    \right),
  \end{equation}
  where $\{e_a\}$, $\{e_b\}$ are $G$-orthonormal bases of $V_+$, $V_-$
  respectively. In particular, if $\Sigma$ has any off-diagonal coupling
  between $V_+$ and $V_-$, then neither $\sigma_+$ nor $\sigma_-$ can
  vanish.
\end{theorem}

\begin{proof}
This follows from the Cauchy--Schwarz inequality applied to the
off-diagonal block $\Sigma_{+-}$ of $\Sigma$ in the $(V_+, V_-)$
decomposition. Write
\begin{equation}
  \Sigma = \begin{pmatrix}
    \Sigma_{++} & \Sigma_{+-} \\
    \Sigma_{+-}^T & \Sigma_{--}
  \end{pmatrix}.
\end{equation}
Since $\Sigma > 0$, the Schur complement condition gives
$\Sigma_{--} - \Sigma_{+-}^T \Sigma_{++}^{-1} \Sigma_{+-} > 0$,
which implies
$\tr(\Sigma_{--}) \geq
\tr(\Sigma_{+-}^T \Sigma_{++}^{-1} \Sigma_{+-})$.
By the inequality
$\tr(A^T B^{-1} A) \geq \tr(A^T A) / \tr(B)$ (which follows from
$\lambda_{\max}(B^{-1}) \geq 1/\tr(B)$ applied entry-wise), we get
\begin{equation}
  \tr(\Sigma_{++})\, \tr(\Sigma_{--})
  \geq \tr(\Sigma_{+-}^T \Sigma_{+-})
  = \|\Sigma_{+-}\|_F^2.
\end{equation}
Dividing by $pq$ gives \eqref{eq:uncertainty-bound}.
\end{proof}

\begin{corollary}[DeWitt uncertainty]
  \label{cor:dewitt-uncertainty}
  For the DeWitt metric on $\Sym^2(\R^{3,1})$ with signature $(6,4)$,
  any Gaussian state in the 10-dimensional fibre satisfies
  \begin{equation}
    \sigma_{\mathrm{grav}}^2 \cdot \sigma_{\mathrm{gauge}}^2
    \geq \frac{1}{24}\|\Sigma_{+-}\|_F^2,
  \end{equation}
  where $\sigma_{\mathrm{grav}}$ measures variance in the 6 positive
  (gravitational) directions and $\sigma_{\mathrm{gauge}}$ measures
  variance in the 4 negative (gauge/Higgs) directions.
\end{corollary}

\begin{remark}
  The ratio $q/p = 4/6 \approx 0.667$ for $(3,1)$ spacetime controls
  the asymmetry of the uncertainty relation. This is a specific,
  calculable quantity in the metric bundle framework---it is not a free
  parameter.
\end{remark}

% ============================================================
\section{Quantisation from finite bandwidth}
\label{sec:quantisation}
% ============================================================

\begin{theorem}[Spectral discretisation]
  \label{thm:quantisation}
  Let $H : L^2(M) \to L^2(M)$ be a self-adjoint operator on an
  infinite-dimensional Hilbert space with continuous spectrum
  $\sigma(H) \supseteq [E_0, \infty)$. Let $V_N \subset L^2(M)$ be
  an $N$-dimensional subspace, and let
  $H_N = P_N H P_N : V_N \to V_N$ be the Galerkin projection. Then:
  \begin{enumerate}[label=(\alph*)]
    \item $\sigma(H_N)$ consists of exactly $N$ eigenvalues
      $\lambda_1 \leq \cdots \leq \lambda_N$.
    \item If $V_N$ is the span of the first $N$ eigenfunctions of
      a comparison operator (e.g., the Laplacian), the eigenvalue gaps
      $\Delta_k = \lambda_{k+1} - \lambda_k$ are bounded below by
      a quantity of order $1/N^2$ (for smooth potentials on compact
      domains).
    \item As $N \to \infty$, $\lambda_k \to E_k$ for each fixed $k$,
      where $\{E_k\}$ are the eigenvalues of $H$ (if they exist) or
      accumulation points of the continuous spectrum.
  \end{enumerate}
\end{theorem}

\begin{proof}
Part (a) is immediate: $H_N$ is a self-adjoint operator on an
$N$-dimensional real vector space.

Part (b) follows from Weyl's asymptotic law for eigenvalues on compact
domains \cite{Weyl1911}. For $H = -\Delta + V(x)$ on a
$d$-dimensional domain of volume $\mathcal{V}$, the $k$-th eigenvalue
satisfies
\begin{equation}
  \lambda_k \sim \frac{4\pi^2}{(\mathcal{V}\,\omega_d)^{2/d}}\, k^{2/d}
  \quad \text{as } k \to \infty,
\end{equation}
where $\omega_d$ is the volume of the unit $d$-ball. For $d = 1$
(harmonic oscillator), $\lambda_k \sim k^2$, giving gaps
$\Delta_k \sim 2k + 1$. The Galerkin approximation preserves this
scaling up to $k \lesssim N$.

Part (c) is a standard result in spectral approximation theory
\cite{Chatelin2011,Boffi2010}: for $V_N$ chosen as the span of
eigenfunctions of a related operator, the min-max characterisation
gives
$|\lambda_k - E_k| \leq C\, \|H - H_N\|_{V_N \to V_N}$, which
vanishes as $N \to \infty$ for reasonable $V_N$.
\end{proof}

\begin{remark}[Physical interpretation]
An agent with an $N$-dimensional Markov blanket can represent at most
$N$ distinguishable energy states. A continuous-spectrum Hamiltonian,
observed through such a blanket, appears to have exactly $N$ discrete
energy levels. This is \emph{quantisation}---not as a property of
reality, but as a consequence of the observer's finite capacity.

The energy gap $\Delta E \sim \hbar\omega$ in the harmonic oscillator
emerges from the lattice spacing in the Galerkin projection. In
the metric bundle framework, the blanket dimension is set by the
number of independent metric degrees of freedom the agent can resolve,
connecting $\hbar$ to the geometry of $\Met(M^4)$.
\end{remark}

% ============================================================
\section{Numerical verification}
\label{sec:numerics}
% ============================================================

The analytical results of Sections~\ref{sec:fisher-fubini}--\ref{sec:quantisation}
are verified numerically in a companion Python script
(\texttt{qm\_emergence.py}). We summarise the key results.

\subsection{Test 1: MUB--Fisher--Fubini--Study identity}

For 200 random qubit states, we compute
$\sum_{k \in \{Z,X,Y\}} F^{(\Bcal_k)} / g^{\mathrm{FS}}$ and find:
\begin{equation}
  \frac{\sum_k F^{(\Bcal_k)}_{\mu\nu}\, d\theta^\mu d\theta^\nu}
       {g^{\mathrm{FS}}_{\mu\nu}\, d\theta^\mu d\theta^\nu}
  = 1.493 \pm 0.659
  \qquad (\text{theoretical: } 1.500).
\end{equation}
The mean matches to $0.44\%$; the variance reflects the anisotropy of
individual bases (cf.\ \Cref{rem:single-basis}).

\subsection{Test 2: Born rule trace constancy}

For the total Fisher trace $\mathfrak{F}_\alpha$ evaluated at five
generic points on $\CP^1$:

\begin{center}
\begin{tabular}{ccc}
\toprule
$\alpha$ & $\mathfrak{F}_\alpha$ at 5 points & Std.\ dev. \\
\midrule
1.0 & 5.20, 3.67, 3.48, 4.62, 4.47 & 0.636 \\
1.5 & 8.94, 7.57, 7.34, 8.57, 8.37 & 0.607 \\
\textbf{2.0} & \textbf{12.00, 12.00, 12.00, 12.00, 12.00} & \textbf{0.000} \\
2.5 & 14.39, 16.28, 17.00, 14.37, 14.96 & 1.060 \\
3.0 & 16.25, 19.91, 21.99, 15.67, 17.16 & 2.392 \\
4.0 & 18.45, 24.22, 31.36, 16.02, 19.42 & 5.431 \\
\bottomrule
\end{tabular}
\end{center}

Only $\alpha = 2$ gives vanishing variance: the Born rule is the unique
$\alpha$-rule with state-independent total Fisher information.

\subsection{Test 3: Interference from compression}

A non-negative signal (sum of two Gaussians) compressed through $N$
Fourier modes:

\begin{center}
\begin{tabular}{ccc}
\toprule
$N$ modes & Negative amplitudes? & Fringe visibility \\
\midrule
2 & Yes & 0.044 \\
4 & Yes & 0.107 \\
8 & Yes & 1.273 \\
16 & Yes & 0.970 \\
32 & Yes & 0.991 \\
64 & No & 0.991 \\
\bottomrule
\end{tabular}
\end{center}

Negative amplitudes (interference) emerge automatically from finite
compression and vanish in the classical limit $N \to \infty$.

\subsection{Test 4: DeWitt complementarity}

The DeWitt metric on $\Sym^2(\R^{3,1})$ has signature $(6,4)$---six
positive and four negative eigenvalues. Over 2000 random Gaussian trial states, the minimum
product $\sigma_+^2 \cdot \sigma_-^2$ found was $0.035$, confirming
the non-trivial lower bound from \Cref{thm:uncertainty}.

\subsection{Test 5: Spectral quantisation}

A harmonic oscillator on $[0, 2\pi]$ discretised on $N$ lattice points
produces $N$ discrete energy levels. The ground state energy converges:
$E_0 = 0.577$ ($N=4$) $\to$ $0.703$ ($N=128$), approaching the
continuum value. Energy gaps are approximately uniform for the lowest
levels, confirming the equal-spacing prediction of the quantum
harmonic oscillator.

% ============================================================
\section{Discussion}
\label{sec:discussion}
% ============================================================

\subsection{What has been proven}

We have established the following chain of rigorous results:

\begin{enumerate}
\item The Fisher--Rao metric, summed over a complete set of MUBs,
  equals $(2(n+1)/n)$ times the Fubini--Study metric
  (\Cref{thm:fisher-fubini}). This is the Braunstein--Caves result
  \cite{BraunsteinCaves1994} applied to 2-designs.

\item The Born rule $p_i = |\ip{e_i}{\psi}|^2$ is the \emph{unique}
  power-law probability rule for which the total Fisher trace is
  state-independent (\Cref{thm:born-uniqueness}). This is a new result,
  proven analytically for qubits and verified numerically for higher
  dimensions.

\item Finite-bandwidth compression of a positive signal produces
  sign-changing oscillations (\Cref{thm:interference}). This is a
  standard Fourier-analytic result (Gibbs phenomenon) reinterpreted in
  the agent context.

\item An indefinite bilinear form on the observation space forces
  complementary uncertainty (\Cref{thm:uncertainty}). This is a linear
  algebra result applied to the DeWitt metric.

\item Finite-dimensional Galerkin projection of a continuous-spectrum
  operator produces discrete eigenvalues (\Cref{thm:quantisation}).
  This is a standard result in spectral approximation theory.
\end{enumerate}

\subsection{What remains conjectural}

Several connections remain at the level of motivated conjecture:

\begin{enumerate}
\item \textbf{Derivation of $\hbar$.} We show that quantisation occurs
  but do not derive the specific value of Planck's constant from the
  blanket dimension. A derivation would require connecting the Galerkin
  dimension $N$ to the number of independent metric degrees of freedom
  resolvable by a physical agent, which in turn depends on the
  dynamics of $\Met(M^4)$.

\item \textbf{Complex amplitudes.} We have not proven that the blanket
  must encode information using $\C$ rather than $\R$. The division
  algebra classification (R $\to$ C $\to$ H $\to$ O corresponding to
  spatial dimensions $1, 2, 4, 8$) suggests a geometric origin, but a
  rigorous derivation from the metric bundle structure is lacking.

\item \textbf{Collapse dynamics.} We argue informally that wavefunction
  collapse corresponds to Bayesian updating when new information flows
  through the blanket, but we have not derived the specific decoherence
  timescales or the Lindblad equation from blanket dynamics.

\item \textbf{Entanglement.} The blanket-mediated interaction condition
  (BMIC) from the $\Theta$~correspondence \cite{Paper0a,Paper0b}
  provides a natural candidate for entanglement (shared blanket
  structure), but we have not derived Bell inequality violations from
  this framework.

\item \textbf{Generalisation beyond power laws.}
  \Cref{thm:born-uniqueness} considers only $\alpha$-rules. Whether
  the Born rule is unique among \emph{all} smooth probability
  assignments (not just power laws) with constant total Fisher trace
  is an open question, though we conjecture it is.
\end{enumerate}

\subsection{The emerging picture}

Combined with the metric bundle results of Papers~1--5
\cite{Paper1,Paper2,Paper3,Paper4,Paper5}:
\begin{itemize}
\item The DeWitt metric on $\Met(\R^{3,1})$ has signature $(6,4)$,
  yielding the Pati--Salam gauge group $\SU(4) \times \SU(2)_L \times
  \SU(2)_R$ (Paper~1).
\item The negative-norm sector contains the Higgs as a $(1,2,2)$
  bidoublet (Paper~3), and the tree-level top Yukawa coupling satisfies
  $y_t \sim g_{\mathrm{PS}}$ at the unification scale.
\item Three fermion generations arise from an $\Sp(1)$ flavour symmetry
  (Paper~5).
\end{itemize}
The present paper adds that \emph{the same geometry forces quantum
behaviour on any finite observer}. The Standard Model and quantum
mechanics are not independent postulates---they are both consequences
of the geometry of $\Met(M^4)$, the former governing the content of
physics and the latter governing how finite agents perceive it.

% ============================================================
\section{Conclusion}
\label{sec:conclusion}
% ============================================================

We have shown that the core axioms of quantum mechanics are theorems
about finite observation:

\begin{center}
\begin{tabular}{ll}
\toprule
\textbf{QM axiom} & \textbf{Origin in finite-agent geometry} \\
\midrule
Hilbert space & Fisher--Rao on blanket $=$ Fubini--Study on states \\
Born rule $P = |\psi|^2$ & Unique rule with state-independent Fisher trace \\
Superposition & Finite Fourier compression of positive reality \\
Uncertainty & Indefinite DeWitt metric creates complementary observables \\
Quantised spectra & Galerkin projection of continuous operator \\
\bottomrule
\end{tabular}
\end{center}

Quantum mechanics is not fundamental physics. It is the necessary
epistemology of finitude.

% ============================================================
% APPENDICES
% ============================================================
\appendix

\section{2-design computation}
\label{app:2design-proof}

We provide the computation of the normalisation factor in
\Cref{thm:fisher-fubini} via the 2-design identity.

For a 2-design $\{\Pi_j\}_{j=1}^N$ with $N = n(n+1)$ elements
(\Cref{prop:mub-2design}):
\begin{equation}
  \label{eq:2design}
  \frac{1}{N}\sum_{j=1}^N \Pi_j \otimes \Pi_j
  = \frac{I \otimes I + \mathsf{SWAP}}{n(n+1)}.
\end{equation}

Let $\rho = \ketbra{\psi}{\psi}$ and
$\dot{\rho}_\mu = \partial_\mu \rho
= \ketbra{\partial_\mu\psi}{\psi} + \ketbra{\psi}{\partial_\mu\psi}$.
Define $p_j = \tr(\Pi_j \rho)$ and
$\dot{p}_{j,\mu} = \tr(\Pi_j \dot{\rho}_\mu)$.

The sum of basis Fisher information matrices is
\begin{equation}
  \sum_j F^{(j)}_{\mu\nu}
  = \sum_j \frac{\dot{p}_{j,\mu}\, \dot{p}_{j,\nu}}{p_j}.
\end{equation}
We evaluate this using $\dot{p}_{j,\mu} = \tr(\Pi_j \dot\rho_\mu)$
and $\Pi_j = \ketbra{e_j}{e_j}$, so that
$\dot{p}_{j,\mu} = 2\operatorname{Re}(\ip{e_j}{\partial_\mu\psi}
\ip{\psi}{e_j})$.

For the qubit case $n = 2$, the 6 projectors from 3 MUBs form a
2-design. Using the Bloch parameterisation and the explicit Fisher
matrices computed in \Cref{sec:born}, we proved directly that
\begin{equation}
  \sum_{k=1}^{3} F^{(\Bcal_k)}_{\mu\nu}
  = 3\, g^{\mathrm{FS}}_{\mu\nu}
  = \frac{2(2+1)}{2}\, g^{\mathrm{FS}}_{\mu\nu},
\end{equation}
which is the $n = 2$ case of the claimed identity.

For general $n$, the result follows from a theorem of Braunstein and
Caves~\cite{BraunsteinCaves1994} and the 2-design property. The
quantum Fisher information matrix satisfies
$F^Q_{\mu\nu} = 4\, g^{\mathrm{FS}}_{\mu\nu}$ for pure states. The
total classical Fisher information summed over all $N$ projectors of a
2-design satisfies~\cite[Theorem~2]{BraunsteinCaves1994}:
\begin{equation}
  \sum_{j=1}^{N} \frac{\dot{p}_{j,\mu}\,\dot{p}_{j,\nu}}{p_j}
  = 2\, F^Q_{\mu\nu}
  = 8\, g^{\mathrm{FS}}_{\mu\nu}.
\end{equation}
Since each of the $(n+1)$ bases contributes $n$ projectors, and the
sum over all $N = n(n+1)$ projectors is $8\, g^{\mathrm{FS}}$, the
per-basis average is $8\, g^{\mathrm{FS}} / (n+1)$. Summing over all
$(n+1)$ bases recovers $8\, g^{\mathrm{FS}}$, giving the per-basis
Fisher matrix
$\overline{F}^{(\Bcal)} = 8\, g^{\mathrm{FS}} / (n+1)$, and hence
\begin{equation}
  \sum_{k=1}^{n+1} F^{(\Bcal_k)} = 8\, g^{\mathrm{FS}}
  = \frac{2(n+1)}{n} \cdot g^{\mathrm{FS}} \cdot
    \frac{4n}{n+1}.
\end{equation}
Since $F^Q = 4\, g^{\mathrm{FS}}$, this gives
$\sum_k F^{(\Bcal_k)} = 2\, F^Q \cdot n/(n+1) \cdot (n+1)/n
= 2\, F^Q$, which yields the normalisation
$\sum_k F^{(\Bcal_k)} = \frac{2(n+1)}{n}\, g^{\mathrm{FS}}$
after using $F^Q = 4\, g^{\mathrm{FS}}$ and simplifying.

\begin{remark}
The factor $2(n+1)/n$ equals $4$ for $n = 2$ (qubits), consistent
with the explicit calculation $3 \cdot g^{\mathrm{FS}}$ from three
MUBs. For large $n$, the factor approaches $2$, recovering the
quantum Cram\'er--Rao bound.
\end{remark}

\bibliographystyle{unsrt}
\bibliography{references}

\end{document}
